\documentclass[10pt,a4paper]{amsart}
\usepackage[margin=19mm]{geometry}
\usepackage{amsmath,amssymb,mathtools}
\usepackage[scr=rsfs,cal=euler]{mathalfa}
\usepackage{xfrac}
\usepackage[unicode,hidelinks,bookmarks=false]{hyperref}
\newcommand{\ie}{i.e.\ }
\newcommand{\cf}{cf.\ }
\newcommand{\resp}{respectively\ }
\newcommand{\Subsection}{\subsection}
\providecommand{\nobreakdash}{\nobreak\hbox{-}}
\setlength{\emergencystretch}{2em}
\multlinegap0pt
\usepackage{csquotes}
\usepackage{mathtools}
\newtheorem{lemma}{Lemma}
\newcommand{\R}{\mathbb{R}}
\title{The Follow-The-Leader scheme with non-monotone velocity}
\author{Marco Di Francesco}
\date{Source: arXiv:2607.28195v1 (CC BY 4.0)}
\begin{document}
\maketitle

\noindent\emph{Source and licence.} The Follow-The-Leader scheme with non-monotone velocity. Version arXiv:2607.28195v1 (CC BY 4.0), distributed by arXiv under the Creative Commons Attribution 4.0 International licence, http://creativecommons.org/licenses/by/4.0/, as recorded in the arXiv licence metadata for this exact version and shown on the abstract page https://arxiv.org/abs/2607.28195v1. Full licence notice: \texttt{licenses/arxiv-2607.28195-CC-BY-4.0.txt}.\par\smallskip

\noindent\textit{Editorial scope.} Counted unit: the article's lemma lem:max\_min (source0.tex lines 143--150). The article's complete original proof of it is delivered with it (source0.tex lines 562--568, one \texttt{proof} environment of 7 lines, which the article places away from the statement). No mathematics is written, repaired or re-derived: the statement and the proof are the article's own lines, in the article's own notation, and no retained line is substituted or re-flowed.  No further source-local context is needed: every cross-reference of every retained line resolves inside the two delivered spans.  Nothing was omitted from any retained span, so the package carries no omission record.  No retained line contains a citation, so the wrapper carries no bibliography and no bibliography entry.  No retained line contains a figure environment, an image-inclusion command or drawing code, so no image is delivered and the package declares no asset dependency.  Editorial formatting only, in addition: an AMS \texttt{amsart} preamble that restates the article's own declarations and macros used by the retained lines, namely the environments \texttt{lemma} on separate counters, together with the standard packages those lines need.  The retained environments print in this document as Lemma 1 in order of appearance, whereas the article prints the same items with the numbering of its own full document.  The complete native source of the article as distributed in the arXiv e-print for this exact version is bundled unchanged as \texttt{sources/206397/source0.tex}.
\begin{lemma}\label{lem:max_min}
    Assume $v$ satisfies (Vel). Let $F:[0,+\infty)^2\rightarrow \R$ be defined by 
    \[F((a,b))=\begin{cases}
        \max_{c\in [b,a]}v(c) & \hbox{if $b\leq a$}\\
        \min_{c\in [a,b]}v(c) & \hbox{if $a<b$}
    \end{cases}\]
    Then, $F$ is a locally Lipschitz function.
\end{lemma}

\begin{proof}[Proof of Lemma \ref{lem:max_min}]
On the domain $D=\left\{(a,b)\in \R^2\,:\,\, a\leq b\right\}$, the map $F$ is trivially Lipschitz continuous with respect to $a$ because the map $[0,b]\ni a \mapsto F(a,b)$ coincides in this case with the non-decreasing hull of $v$ on $[0,b]$. Similarly, the map $[a,+\infty)\ni b \mapsto F(a,b)$ being the non-decreasing hull of $v$ on $[a,+\infty)$ implies the Lipschitz continuity with respect to $b$. A similar argument shows the (joint) Lipschitz continuity of $F$ on $\R^2\setminus D$. The only case left to check the Lipschitz continuity with respect to $a$ is the one in which $(a_1,b)\in D$ and $(a_2,b)\not\in D$. In that case, for some $\hat{a}_1\in [a_1,b]$ and for some $\hat{a}_2\in [b,a_2]$, we have
\begin{align*}
    & |F(a_2,b)-F(a_1,b)|=|v(\hat{a}_2)-v(\hat{a}_1)|\leq [v]_{\mathrm{Lip}[a_1,a_2]}|\hat{a}_1-\hat{a}_2|\leq [v]_{\mathrm{Lip}[a_1,a_2]}|a_1-a_2|\,.
\end{align*}
A similar estimate applies to the case $b_1\leq a\leq b_2$, which completes the proof.
\end{proof}

\end{document}
